\subsection{Tanzimatın Uygulaması ve Karşılaşılan Zorluklar}
\indent\indent Daha önce de belirtildiği gibi, Tanzimat Fermanı'na kadar devletin altın ya da klasik çağı olan Kanuni dönemine dönüşün yolları kânun-ı kadîm mefhumu üzerinden ordu temelli reformlarda aranmıştır.
Tanzimat ile birlikte, Kanuni dönemine dönüşün artık mümkün olmadığı kabul edilmiştir. Çünkü Kanuni'nin altın çağı iki yüz sene öncesinde kalmıştır. Dünya, iki yüz sene önceki dünya değildir. 
Kalemiye sınıfı, yeni altın çağın inşasına giden yolun batı tarzı modern müesseselerden geçtiğini düşünmektedir. Bürokrasideki bu radikal reform, birçok kesim tarafından tepki topladı. Oluşan tepkilerin 
detaylı analizi bu araştırmanın amacı dışındadır. Yine de Tanzimat'a karşı oluşan ana muhalefet noktalarını belirtmek faydalı olacaktır.

\subsection{Kalemiye-İlmiye Çatışması}
\indent\indent Tanzimat ile birlikte, kalemiye-ilmiye çatışması ortaya çıkmıştır. Tanzimat'a kadar taşra idaresi, kadının başkanlığında toplanan mahalli meclisler ile sağlanıyordu. 
Tanzimat ise başkanlığı kadıdan(ilmiyeden) alıp vali, muhassıl, kaza müdürü gibi kalemiye memurlarına vermiştir. Bununla birlikte, gayrimüslim teb`anın dini liderleri ve kocabaşıları da bu meclislere
üye olarak dahil olabiliyorlardı. İdare adamlarının başkanlığı ile taşra idaresinin ulemanın ve ayânın nüfuzundan kurtarılması amaçlanmıştı.\footcite[181]{tanzimat}

Diğer bir tartışma ise ceza hukuku alanında yaşanmıştır. 1840 Ceza Kanunnamesi, kanunnamenin ihtiva ettiği suçların icra-ı mahkeme yetkisini mahalli meclislere ve Meclis-i Ahkâm-ı Adliye'ye vermiştir. Bu noktada
kanunnamenin ikinci faslındaki üçüncü ve dördüncü maddelere göz atmak faydalı olacaktır:
\begin{itemize}
    \item \textbf{Üçüncü Madde}\\
    Bu misillü fesad-ı kavli ve fi'liye cesaret edenler Dersaadet'de oldukları halde davalarının rü'yeti ve tevâtüren(kesin olarak isbatı) mutlaka Meclis-i Ahkâm-ı Adliyeye mahsus ve münhasır buluna.
    \item \textbf{Dördüncü Madde}\\
    Taşralarda vuku bulduğu takdirde dahi evvel emirde memleket meclisinde görülüp şahitler Dersaadet'e celp ile kezalik Meclis-i Ahkâm-ı Adliyede tekrar davası görüle.\footcite[812]{islam_osmanli_hukuku}
\end{itemize}

Çatışmanın ana sebebi, Osmanlı bürokrasisindaki iki grubun iktidar mücadelesi olduğunu söylersek abartmış olmayız. Kalemiye sınıfı, seküler müesseseler ile ilmiye sınıfının siyaset alanını daratlmayı amaçlıyordu. 
Eski müesseseleri savunan muhafazakar ilmiye sınıfının bu durumu öylece kabul etmesi beklenemezdi. Nitekim, Reşit Paşa'nın 31 Mart 1841'deki azli sonrasında muhafazakarlar iktidara gelmiştir. Yeni idare askeri ıslahata ve 
müslümanları tatminiyete öncelik vermiştir.\footcite[185]{tanzimat}

\subsection{Vergi Toplamadaki Reformlar}
\indent\indent Tanzimat'ın getirdiği en radikal reform iltizamın kaldırılması olmuştur. Meclis-i Vâlâ iltizam yerine servete dayalı yeni bir vergilendirme sistemini getiriyordu. Yeni vergi sistemine göre,
yüz akçe servetten üç akçe on iki paralık bir vergi alınacaktı.\footnote{Bir akçe kırk paraya denk gelmekteydi.} Verginin toplamasında mültezim devre dışı kalıyor; yetki, yukarıda değinilen merkezden 
atanan memurların başkanlığında toplanan  meclislere veriliyordu. Merkezi yönetim, vergi kaybını ve istismarı en aza indirmek amacıyla kontrolü ele almayı amaçlıyordu.

Tahmin edileceği üzere bu uygulama mültezim, voyvoda gibi bu yolla zenginleşen geniş bir zümrenin gelir ve istismar kapısını kapatıyordu. Öte yandan vergi ödemede eşitliğin sağlanarak
imtiyaz ve muafiyetlerin kaldırılması, eskiden az vergiveren âyanları, ağaları ve din adamlarını ıslahat aleyhine çevirmiştir.\footcite[180]{tanzimat}

Oluşan muhalefet ve yeterli memurun bulunmaması sebebiyle 1839 ve 1840 senelerinde vergi gelirlerinin önemli bir kısmı toplanamamış, bütçe açık vermiştir. Reşit Paşa, devlet masraflarını ve 
maaşları ödemek için istikrazı denemiş ve \textit{eshâm kavâ} denilen bonolarla ödeme yapmıştır. Mehmet Ali Paşa ile olan mücadeleden dolayı zaten sıkıntılı olan hazine artık krize girmiştir.
Bu noktadan sonra Reşit Paşa'nın azli kaçınılmaz hale gelmişti.\footcite[185]{tanzimat}
